Cieľom tejto diplomovej práce bolo navrhnúť a~implementovať
systém na rozpoznávanie gest rozhodcu zápasenia pomocou
inteligentných hodiniek Amazfit T-Rex 3 s~operačným systémom
Zepp OS a~sprievodnej mobilnej aplikácie pre Android. Práca
prepája viaceré oblasti -- šport, pohybové senzory
v~inteligentných hodinkách a~vývoj softvéru pre dve odlišné
platformy -- a~ústi do funkčného riešenia, ktoré dokáže
v~reálnom čase zaznamenávať udalosti v~zápase na základe
pohybu zápästia rozhodcu.

V~teoretickej časti sme sa najprv venovali základným pojmom
z~oblasti inteligentných hodiniek, operačnému systému Zepp~OS
a~zabudovaných senzorov. Osobitnú pozornosť sme venovali
akcelerometru a~gyroskopu, ktoré tvoria základ rozpoznávania
gest. Následne sme analyzovali zápasenie ako šport -- jeho
hlavné štýly, priebeh zápasu, úlohy členov rozhodcovského
tímu a~situácie, v~ktorých rozhodca udeľuje body, napomenutia
alebo signalizuje pasivitu. Na základe pravidiel organizácie
\Gls{uww} sme identifikovali kľúčové signály, ktoré sú pre
priebeh zápasu nevyhnutné, a~ktoré sa zároveň dajú zachytiť
pomocou senzorov v~hodinkách.

V~ďalšej časti sme predstavili vývojové prostriedky použité
pri tvorbe systému -- JavaScript a~Zeus \Gls{cli} pre vývoj
aplikácie pre Zepp~OS, Kotlin a~Android Studio pre mobilnú
aplikáciu a~Python s~knižnicami Pandas a~Matplotlib na
spracovanie a~vizualizáciu nameraných dát. Po definovaní
požiadaviek na hardvér -- predovšetkým prístupu k~surovým
dátam zo senzorov a~podpory vlastných aplikácií -- sme zvolili
hodinky Amazfit T-Rex~3, ktoré tieto požiadavky spĺňajú
a~zároveň poskytujú dostatočný výkon pre snímanie pohybu
v~reálnom čase.

Analýza existujúcich riešení potvrdila, že hoci v~športovom
prostredí existuje viacero rozhodcovských aplikácií pre smart
hodinky -- napríklad pre futbal alebo hokej -- špecifická
aplikácia pre rozhodcov v~zápasení s~detekciou gest zatiaľ
chýba. Aplikácie zamerané na zápasenie sú obmedzené na ručný
zápis bodov prostredníctvom dotykového displeja, čo nevyužíva
plný potenciál inteligentných hodiniek. Existujúce akademické
práce v~oblasti rozpoznávania gest pomocou inerciálnych
senzorov nám zároveň poskytli odrazový mostík pre vlastný
návrh.

Jadrom práce bol návrh a~implementácia samotných gest. Pre
ovládanie aplikácie sme zaviedli niekoľko elementárnych
pohybov ruky -- zdvihnutie ruky nad hlavu, presun ruky za
chrbát, rotáciu zápästia a~návrat ruky do neutrálnej polohy
-- z~ktorých sa skladajú kompozitné gestá zodpovedajúce
signálom rozhodcu: BODOVANIE, NAPOMENUTIE, PASIVITA
a~TOUCHE, s~rozlíšením medzi červeným a~modrým zápasníkom.
Návrh gest sme zámerne držali blízko bežnej gestikulácie
rozhodcu na žinenke, aby si rozhodca nemusel osvojovať umelé
pohyby a~aby gestá zostali viditeľné pre ostatných členov
rozhodcovského tímu.

Samotná implementácia pozostáva z~dvoch častí -- aplikácie
pre Zepp~OS na hodinkách, ktorá zbiera dáta z~akcelerometra
a~gyroskopu a~odosiela ich cez \Gls{ble} do Android
aplikácie, a~mobilnej Android aplikácie, kde prebieha
vlastné rozpoznávanie gest na základe nakalibrovaných
prahových hodnôt (\uv{threshold-based} algoritmus). Súčasťou
implementácie je aj vibračná spätná väzba na hodinkách,
ktorá rozhodcovi v~reálnom čase potvrdzuje rozpoznané gestá
bez potreby pozerať na displej. Pre potreby ladenia
a~kalibrácie sme aplikáciu rozšírili o~vývojový režim,
v~ktorom je možné sledovať priebeh detekcie spolu so
surovými dátami zo senzorov.

Testovanie aplikácie sme rozdelili do troch fáz, ktoré sa
postupne približovali k~reálnym podmienkam použitia. V~prvej
fáze sme overili rozpoznávanie izolovaných gest v~ideálnych
podmienkach -- naprieč 160~pokusmi (osem gest po
20~opakovaní) aplikácia dosiahla celkovú úspešnosť 99,4\,\%,
pričom sedem z~ôsmich gest bolo rozpoznaných so 100\,\%
úspešnosťou. V~druhej fáze sme overili spoľahlivosť aplikácie
v~kontexte kompletného zápasníckeho scenára -- 10-násobné
opakovanie 13-krokovej sekvencie viedlo k~presnosti 96,2\,\%
s~nulovou mierou falošných detekcií. V~tretej fáze sme
testovali aplikáciu s~dvomi externými testermi bez
predchádzajúcej skúsenosti so zápasníckou gestikuláciou,
ktorí dosiahli presnosti 83,1\,\% a~90,8\,\%. Celkovo sme
naprieč všetkými fázami vyhodnotili 420~pokusov o~detekciu
gesta.

Z~výsledkov vyplýva, že aplikácia preukázala vysokú
spoľahlivosť pre kalibračného používateľa a~rozumnú
prenositeľnosť na ďalších používateľov bez explicitnej
kalibrácie. Pri testovaní sme identifikovali dve hlavné
obmedzenia systému. Prvým je závislosť prahového algoritmu
na individuálnom štýle vykonávania gest, ktorá vedie
k~poklesu presnosti pri externých používateľoch. Toto
zistenie sa zhoduje s~výsledkami minuloročnej
práce~\cite{dp-balla}, čo potvrdzuje, že ide o~systémový
jav prahového prístupu nezávislý od konkrétnej hardvérovej
platformy. Druhým obmedzením je hardvérový limit samotných
hodiniek pri rýchlych po sebe nasledujúcich pohyboch
zápästia -- pri viacnásobnom udelení bodov v~rýchlom slede
senzor gyroskopu v~niektorých prípadoch zaznamená iba časť
vykonaných pohybov. Toto obmedzenie nie je možné odstrániť
na úrovni našej aplikácie a~v~praxi sa dá zmierniť
ponechaním dostatočného časového odstupu medzi jednotlivými
pohybmi.

Do budúcnosti existuje viacero smerov, ktorými je možné
projekt rozšíriť:

\begin{itemize}
    \item doplnenie kalibračnej fázy pred prvým
        použitím aplikácie, v~ktorej by si rozhodca vykonal
        každé gesto niekoľkokrát a~hraničné hodnoty by sa
        automaticky prispôsobili,
    \item prehodnotenie návrhu gesta \uv{TOUCHE}
        s~ohľadom na anatomické rozdiely medzi
        používateľmi, aby bolo vykonateľné bez ohľadu na
        telesné proporcie,
    \item rozšírenie systému o~druhé hodinky nosené na
        pravom zápästí, keďže rozhodca pri rozhodovaní
        používa obidve ruky -- súčasné riešenie sníma
        pohyb iba z~ľavej ruky, čo obmedzuje priestor
        pre návrh ďalších gest,
    \item rozšírenie implementácie o~ďalšie signály
        rozhodcu, ktoré by využili aj pohyb pravej ruky
        a~doplnili tak chýbajúce udalosti v~priebehu zápasu.
\end{itemize}

Najväčším prínosom tejto práce je preukázanie, že aj bežne
dostupné komerčné inteligentné hodinky -- v~kombinácii so
sprievodnou mobilnou aplikáciou -- môžu plniť úlohu pomocníka
rozhodcu v~zápasení bez potreby špecializovaného hardvéru
či zložitých metód strojového učenia. Navrhnuté riešenie
zachováva prirodzenú gestikuláciu rozhodcu, prispieva
k~pohodlnejšiemu zaznamenávaniu udalostí v~zápase a~zároveň
otvára priestor pre podobné aplikácie v~iných športoch
s~ustálenou gestikulačnou komunikáciou.